%% mwe_fcolorbox_phase3.tex
%% Pergunta: \fcolorbox sobrevive ao tagging phase-III?
%%   (A) duas \fcolorbox{\parbox} ficam LADO A LADO, ou o patch de tagging
%%       força \par em volta e as empilha (como acontece com tcolorbox)?
%%   (B) a árvore de tags fecha sem "text-unit para hooks differ"?
%%   (C) uma \fcolorbox de largura da mancha com fundo tingido serve de
%%       box de destaque?
%%
%% A estrutura de tagging é transplante fiel de colunas/\coluna
%% (guia-visual.sty §9), cuja ordem é:
%%   abertura: \par, \topskip=0pt, paraoff, tagbegin{Sect} container
%%   cada box: tagbegin{Sect}, \parbox{ paraon #corpo\par paraoff }, tagend
%%   fecho:    \par, tagend, paraon
%% Compilar: sh lualatex-podman.sh mwe_fcolorbox_phase3.tex

\DocumentMetadata{
  lang        = pt-BR,
  pdfversion  = 2.0,
  pdfstandard = ua-2,
  testphase   = {phase-III, table, firstaid},
}

\documentclass[11pt]{article}

\usepackage[T1]{fontenc}
\usepackage[a4paper, margin=25mm]{geometry}
\usepackage{xcolor}

% Paleta CRP-SP
\definecolor{herhighness}{HTML}{422C73}
\definecolor{piquantpink}{HTML}{EE05F2}
\definecolor{royalpretender}{HTML}{A55DA6}
\definecolor{blackwash}{HTML}{0D0D0D}
\colorlet{alarmefundo}{piquantpink!10!white}

\setlength{\parindent}{0pt}

% Helpers de tagging com guard (padrão de guia-visual.sty)
\newcommand{\tagbegin}[1]{\ifdefined\tagstructbegin\tagstructbegin{tag=#1}\fi}
\newcommand{\tagend}{\ifdefined\tagstructend\tagstructend\fi}
\newcommand{\paraoff}{\ifdefined\tagpdfparaOff\tagpdfparaOff\fi}
\newcommand{\paraon}{\ifdefined\tagpdfparaOn\tagpdfparaOn\fi}

\newlength{\cardgap}   \setlength{\cardgap}{6mm}
\newlength{\cardwd}
\newlength{\cardparwd}
\newlength{\parskipsalvo}
\AddToHook{begindocument/end}{\setlength{\parskipsalvo}{\baselineskip}}

% \boxlado[corda borda]{rótulo}{título}{corpo}
\newcommand{\boxlado}[5][royalpretender]{%
  \setlength{\fboxrule}{#2}%
  \setlength{\cardparwd}{\dimexpr\cardwd-2\fboxsep-2\fboxrule\relax}%
  \tagbegin{Sect}%
  \fcolorbox{#1}{white}{%
    \parbox[t]{\cardparwd}{%
      \vspace{0pt}%
      \setlength{\parskip}{\parskipsalvo}%
      \paraon
      {\footnotesize\bfseries\color{#1}#3}\par
      {\large\bfseries\color{herhighness}#4}\par
      #5\par
      \paraoff
    }%
  }%
  \tagend
}

\begin{document}

\section*{Teste A --- dois \texttt{fcolorbox} lado a lado}

Parágrafo antes, para dar contexto de fluxo ao tagging.

\setlength{\cardwd}{\dimexpr(\linewidth-\cardgap)/2\relax}
\setlength{\fboxsep}{5mm}

\par
\topskip=0pt\relax
\paraoff
\tagbegin{Sect}%
\noindent
\boxlado[royalpretender]{0.4pt}{OPÇÃO A}{Manter o modelo atual}{%
  Texto do primeiro card. Precisa quebrar em mais de uma linha para
  valer como teste de \texttt{parbox} dentro de \texttt{fcolorbox}.}%
\hspace{\cardgap}%
\boxlado[piquantpink]{1.2pt}{OPÇÃO B $\cdot$ RECOMENDADA}{Nova divisão de trabalho}{%
  Texto do segundo card, também com mais de uma linha, para conferir
  se os topos das duas caixas alinham.}%
\par
\tagend
\paraon

\vspace{6mm}
Parágrafo depois, para confirmar que o fluxo normal foi retomado.

\section*{Teste B --- box de destaque de largura total}

\setlength{\fboxrule}{0.8pt}%
\setlength{\fboxsep}{5mm}%

\par
\paraoff
\tagbegin{Sect}%
\noindent\fcolorbox{piquantpink}{alarmefundo}{%
  \parbox{\dimexpr\linewidth-2\fboxsep-2\fboxrule\relax}{%
    \setlength{\parskip}{\parskipsalvo}%
    \paraon
    {\footnotesize\bfseries\color{herhighness} CASO DOCUMENTADO $\cdot$ 2026}\par
    {\color{blackwash}Em 2026 --- ano com a menor equipe nominal da série,
    2 pessoas --- a edição 209 foi impressa e postada com dois erros
    materiais no expediente.}\par
    \textbf{Nenhuma etapa do processo era capaz de interceptar o erro
    antes da impressão.}\par
    \paraoff
  }%
}%
\par
\tagend
\paraon

\vspace{4mm}
Fim do teste.

\end{document}
